\section{Focused cyclic sequent calculus}
\label{sec:focused}

This section presents a cyclic sequent calculus for the term language of Section~\ref{sec:calculus}. The calculus is designed to mirror supercompilation operations: \emph{driving} (asynchronous unfolding), \emph{splitting} (synchronous case analysis), and \emph{folding} (cyclic backlinks).

\subsection{Sequents and judgements}

A sequent has the form $\Gamma \vdash t : A$, representing the goal of finding a term $t$ of type $A$ in context $\Gamma$. The sequent calculus manipulates these goals using inference rules.

We use three kinds of judgements:
\begin{itemize}
\item \textbf{Typing judgements} $\jTy(\Gamma, t, A)$: prove that $t : A$ in context $\Gamma$
\item \textbf{Observation judgements} $\mathsf{jObs}(\Gamma, t, \phi)$: prove that term $t$ satisfies observation $\phi$
\item \textbf{Substitution judgements} $\jSub(\Delta, v, \Gamma)$: witness that vertex $v$ can substitute for context $\Delta$ in $\Gamma$
\end{itemize}

\subsection{Proof graphs}

A \emph{cyclic proof} is a finite directed graph where:
\begin{itemize}
\item Each vertex $v$ is labeled with a sequent $\labelof(v)$
\item Each vertex has a list of successor vertices $\succof(v)$ (premises)
\item Each vertex has an attached local rule $\Rule(v)$ validating the step
\item The graph may contain cycles (back-edges)
\end{itemize}

\noindent A proof graph is \emph{valid} if:
\begin{enumerate}
\item Every vertex's local rule is valid (correct premises for the conclusion)
\item Every infinite path contains a \emph{progress event} (a synchronous split, see Section~\ref{sec:global-condition})
\end{enumerate}

\subsection{Asynchronous rules (driving)}

These rules correspond to CBN reduction steps and type-directed decomposition. They are \emph{deterministic} and \emph{invertible}: there is no choice in applying them.

\subsubsection{$\beta$-reduction}

\begin{center}
\AxiomC{$\Gamma \vdash t[u/] : A$}
\UnaryInfC{$\Gamma \vdash \App{(\Lam{B}{t})}{u} : A$}
\DisplayProof
\quad (\textsc{drive-beta})
\end{center}

\subsubsection{Fixpoint unfolding}

\begin{center}
\AxiomC{$\Gamma \vdash t[\Fix{A}{t}/] : A$}
\UnaryInfC{$\Gamma \vdash \Fix{A}{t} : A$}
\DisplayProof
\quad (\textsc{drive-fix})
\end{center}

\subsubsection{Iota reduction (case-of-constructor)}

\begin{center}
\AxiomC{$\Gamma \vdash b_c\,\vec{a} : C[\Roll{I}{c}{\vec{a}}/]$}
\UnaryInfC{$\Gamma \vdash \Case{I}{(\Roll{I}{c}{\vec{a}})}{C}{\vec{b}} : C[\Roll{I}{c}{\vec{a}}/]$}
\DisplayProof
\quad (\textsc{drive-iota})
\end{center}

\subsubsection{Commuting conversions}

When the scrutinee of a case is not a constructor, we attempt to drive it:

\begin{center}
\AxiomC{$\Gamma \vdash t' : \Ind{I}{\vec{a}}$}
\AxiomC{$\Gamma \vdash \Case{I}{t'}{C}{\vec{b}} : C[t'/]$}
\BinaryInfC{$\Gamma \vdash \Case{I}{t}{C}{\vec{b}} : C[t/]$}
\DisplayProof
\quad (\textsc{drive-case-scrut})
\end{center}

\noindent where $t \to t'$ is a CBN reduction step.

\subsection{Synchronous rules (splitting)}

These rules introduce \emph{choice}: they analyze a neutral term (variable, stuck case) by case-splitting on its constructors. This is the only source of nondeterminism in proof search.

\subsubsection{Case split on neutral variable}

\begin{center}
\AxiomC{$\Gamma \vdash b_0 : C[\Roll{I}{0}{}/]$}
\AxiomC{$\cdots$}
\AxiomC{$\Gamma \vdash b_{n-1} : C[\Roll{I}{n-1}{}/]$}
\TrinaryInfC{$\Gamma, x : \Ind{I}{\vec{a}} \vdash \Case{I}{x}{C}{\vec{b}} : C[x/]$}
\DisplayProof
\quad (\textsc{split-case-var})
\end{center}

\noindent Each premise corresponds to one constructor of the inductive type. The context in each premise is extended with the constructor arguments.

\textbf{Key property:} This is a \emph{progress event}. In the cyclic proof, every cycle must pass through at least one split rule. This ensures termination of the proof search (see Section~\ref{sec:global-condition}).

\subsection{Observation rules}

Observation rules prove that a term satisfies a given observation predicate. For natural numbers, the observations are ``reduces to zero'' and ``reduces to successor''.

\subsubsection{Observe zero}

\begin{center}
\AxiomC{}
\UnaryInfC{$\Gamma \vdash \Roll{\Nat}{\mathsf{zero}}{} : \mathsf{Obs}_{\mathsf{zero}}$}
\DisplayProof
\quad (\textsc{obs-zero})
\end{center}

\subsubsection{Observe successor}

\begin{center}
\AxiomC{$\Gamma \vdash n : \Nat$}
\UnaryInfC{$\Gamma \vdash \Roll{\Nat}{\mathsf{succ}}{n} : \mathsf{Obs}_{\mathsf{succ}}$}
\DisplayProof
\quad (\textsc{obs-succ})
\end{center}

\subsubsection{Drive to observation}

\begin{center}
\AxiomC{$\Gamma \vdash t' : \phi$}
\UnaryInfC{$\Gamma \vdash t : \phi$}
\DisplayProof
\quad (\textsc{obs-drive})
\end{center}

\noindent where $t \to t'$ is a CBN reduction step. This allows driving a term until it reaches an observable form.

\subsection{Backlink rules (folding)}

A backlink closes a proof goal by pointing to a previous vertex with a compatible sequent. This is the cyclic aspect of the calculus, and it embodies a fundamental difference from traditional inductive proofs: \emph{the induction principle emerges from the cycle structure rather than being prescribed upfront}.

In traditional induction, one must choose an induction scheme before beginning the proof (e.g., ``induct on list $l$''). In cyclic proofs, the backlink itself \emph{is} the induction hypothesis, discovered during proof search when we recognize that a current goal is an instance of a previous one. The cycle formed by the backlink implicitly defines a recursive proof structure, and the trace condition ensures this recursion is well-founded.

\subsubsection{Exact match}

\begin{center}
\AxiomC{}
\UnaryInfC{$\Gamma \vdash t : A$}
\DisplayProof
\quad (\textsc{back-exact})
\end{center}

\noindent If there exists a prior vertex $v'$ with label $\Gamma \vdash t : A$ (syntactically identical), we can immediately close the goal with a backlink to $v'$.

\subsubsection{Instance match}

\begin{center}
\AxiomC{$\jSub(\Delta, v', \Gamma)$}
\UnaryInfC{$\Gamma \vdash t[\sigma/] : A[\sigma/]$}
\DisplayProof
\quad (\textsc{back-inst})
\end{center}

\noindent If there exists a prior vertex $v'$ with label $\Delta \vdash t : A$ and a substitution $\sigma : \Delta \to \Gamma$, we can close the current goal. The substitution judgement $\jSub(\Delta, v', \Gamma)$ witnesses that $\sigma$ is a valid instantiation.

\textbf{Discovering induction post-hoc:} The instance match rule is particularly powerful because it allows the proof to discover the appropriate instantiation during search. For example:
\begin{itemize}
\item Starting goal: $l : \mathsf{List} \vdash \mathsf{length}\,(\mathsf{map}\,f\,l) : \Nat$
\item After driving and splitting on $l$, in the \texttt{cons} branch we reach: $xs : \mathsf{List} \vdash \mathsf{length}\,(\mathsf{map}\,f\,xs) : \Nat$
\item This matches the original goal with substitution $l \mapsto xs$
\item The backlink implicitly defines structural induction on $l$, discovered post-hoc
\end{itemize}

This flexibility allows cyclic proofs to find induction schemes that would be difficult to specify in advance, especially for mutual recursion, nested induction, or induction on complex measures.

\subsection{Phases and focusing}

We partition the rules into two phases:
\begin{itemize}
\item \textbf{Asynchronous phase}: Apply driving rules eagerly (no backtracking needed)
\item \textbf{Synchronous phase}: Apply split or backlink (requires choice)
\end{itemize}

\noindent This discipline ensures that:
\begin{enumerate}
\item Driving (asynchronous) is always deterministic and eager
\item Splitting (synchronous) is the only choice point
\item Backlinks can be tried at any synchronous phase
\end{enumerate}

\noindent This matches the supercompilation strategy exactly: drive terms until stuck, then split on neutrals or fold via memo table.

\subsection{Global progress condition}
\label{sec:global-condition}

A cyclic proof is \emph{sound} if every infinite path through the graph contains infinitely many progress events. In our setting:
\begin{itemize}
\item A \textbf{progress event} is a synchronous split rule (\textsc{split-case-var})
\item This ensures structural descent: each split analyzes a smaller inductive component
\end{itemize}

\noindent Formally, we attach a \emph{trace state} $\tau$ to each vertex, tracking which inductive variable is being analyzed:
\[
\tau ::= \mathsf{None} \mid \mathsf{Some}(I, x)
\]

\noindent The trace updates as follows:
\begin{itemize}
\item Driving rules: $\tau' = \tau$ (no change)
\item Split rule on variable $x$ of inductive $I$: $\tau' = \mathsf{Some}(I, x)$
\item Backlink: $\tau' = \tau$ (inherit from target)
\end{itemize}

\noindent We define a well-founded order on trace states:
\[
\mathsf{Some}(I, x') <_{\tau} \mathsf{Some}(I, x) \iff x' \text{ is a recursive argument of } x
\]

\noindent A cycle is valid if the trace state strictly decreases on at least one back-edge. This is equivalent to requiring a split event in the cycle, because splits introduce recursive subterms.

\subsection{Relation to supercompilation}

The sequent calculus rules directly correspond to supercompilation operations:
\begin{itemize}
\item \textbf{Driving rules} $\leftrightarrow$ \texttt{drive\_cbn\_once} (CBN unfolding)
\item \textbf{Split rule} $\leftrightarrow$ \texttt{split\_case\_var} (case analysis on neutral)
\item \textbf{Backlink} $\leftrightarrow$ memo table lookup (folding)
\item \textbf{Progress condition} $\leftrightarrow$ termination via structural descent
\end{itemize}

\noindent The bisimulation theorem (Section~\ref{sec:equiv}) makes this correspondence precise: every supercompilation step corresponds to exactly one sequent rule application, and vice versa.

\subsection{Example: \texttt{length (map f l)}}

Consider the fusion of \texttt{length (map f l)}. The sequent proof proceeds as follows:

\begin{enumerate}
\item Start: $l : \mathsf{List} \vdash \mathsf{length}\,(\mathsf{map}\,f\,l) : \Nat$
\item Drive (unfold map): $\vdash \mathsf{length}\,(\Case{\mathsf{List}}{l}{\ldots}{\ldots}) : \Nat$
\item Drive (commute case): $\vdash \Case{\mathsf{List}}{l}{\ldots}{\ldots} : \Nat$
\item Split on $l$:
   \begin{itemize}
   \item Branch \texttt{nil}: $\vdash \mathsf{length}\,\mathsf{nil} : \Nat$ (drive to zero)
   \item Branch \texttt{cons}: $x : A, xs : \mathsf{List} \vdash \mathsf{length}\,(\mathsf{map}\,f\,(\mathsf{cons}\,x\,xs)) : \Nat$
   \end{itemize}
\item In cons branch, drive produces: $\vdash \mathsf{succ}\,(\mathsf{length}\,(\mathsf{map}\,f\,xs)) : \Nat$
\item Backlink to step 1 (with substitution $l \mapsto xs$)
\end{enumerate}

\noindent The cycle passes through the split at step 4, which is a progress event. The trace state tracks descent on $l$ (from $l$ to $xs$), validating the cycle.

\subsection{Vertex classification and correspondence}
\label{sec:vertex-classification}

The connection between the supercompiler's configuration graph and the focused sequent calculus is established by classifying each vertex according to its successor structure. For a vertex $v$ with label $\jTy(\Gamma, t, A)$ and successors $\succof(v) = [w_1, \ldots, w_n]$:

\begin{center}
\small
\begin{tabular}{c|c|c|c}
\textbf{Successors} & \textbf{Rule} & \textbf{SC operation} & \textbf{Phase} \\
\hline
$n = 0$ & \textsc{dr-leaf} & leaf (axiom) & — \\
$n = 1$, $\mathsf{drive\_cbn\_once}(t) \neq t$ & \textsc{dr-cbn-once} & driving (async) & invertible \\
$n = 1$, $\mathsf{drive\_cbn\_once}(t) = t$ & \textsc{dr-fold} & folding (memo hit) & cyclic closure \\
$n > 1$ & \textsc{dr-split} & case-split (sync) & choice point \\
\end{tabular}
\end{center}

\noindent The focused discipline emerges naturally: asynchronous (driving) rules are deterministic and can be applied eagerly; synchronous (split) rules introduce nondeterminism at case boundaries; fold rules close cycles via backlinks to previously seen vertices.

\begin{proposition}[Vertex classification --- mechanized]
\label{prop:vertex-class}
For every vertex $v$ in the configuration graph produced by successful supercompilation, if $\mathsf{cb\_label}(v) = \jTy(\Gamma, t, A)$ and $\mathsf{cb\_succ}(v) = [w_1,\ldots,w_n]$, then the vertex satisfies exactly one of the four rules above. This is proved in Coq by $\texttt{cfg\_vertex\_drive\_rule}$ (\texttt{SupercompilationCorrespondence.v}).
\end{proposition}
